\documentclass[conference]{IEEEtran}

\usepackage{cite}
\usepackage{amsmath,amssymb}
\usepackage{graphicx}
\usepackage{booktabs}
\usepackage{multirow}
\usepackage{array}
\usepackage{xcolor}
\usepackage{url}
\usepackage{makecell}

\title{Class-Balanced Classical Machine Learning for Heart Disease Prediction on the Cleveland Heart Disease Dataset}

\author{
\IEEEauthorblockN{Praveenasarma B.}
\IEEEauthorblockA{
Department of Computer Science \& Engineering\\
University of Moratuwa, Sri Lanka\\
praveenasarma2001@gmail.com
}
}

\begin{document}
\maketitle

\begin{abstract}
This project studies heart disease prediction with classical machine learning on the Cleveland subset of the University of California Irvine (UCI) Heart Disease dataset. The task follows the standard binary formulation in which diagnosis values greater than zero indicate the presence of disease. After reading the supplied task description, dataset notes, benchmark papers, and archive metadata, the experimental file was selected as \texttt{processed.cleveland.data}, because the provided \texttt{WARNING} file states that \texttt{cleveland.data} is corrupted. We first reproduce the benchmark family reported by Nasution \textit{et al.} using Logistic Regression, Random Forest, Support Vector Machine (SVM), and K-Nearest Neighbors (KNN). Baseline models are evaluated with both an 80/20 stratified holdout split and 5-fold stratified cross-validation, using only the training partition for cross-validation so that the final holdout set remains untouched until final evaluation. We then evaluate an additional classical pipeline with feature engineering, class-balanced learning, and hyperparameter tuning, followed by a soft-voting ensemble. On the reproduced holdout split, the strongest raw baseline is KNN with 91.80\% accuracy, 0.9123 F1 score, and 0.9627 receiver operating characteristic area under the curve (ROC-AUC). Across 5-fold stratified cross-validation on the training partition, the strongest baseline by mean F1 score is Logistic Regression without feature selection, with mean accuracy 84.70\% and mean F1 score 0.8254. The evaluated ensemble achieves 88.52\% accuracy, 0.8814 F1 score, and 0.9600 ROC-AUC on the holdout split, but it does not outperform the strongest simple baseline in holdout accuracy or F1 score. Additional ablation analysis shows that class weighting contributes more than feature engineering to cross-validated F1 score in the tuned Logistic Regression family. The study emphasizes transparent reporting, explicit assumptions, and careful interpretation for small medical datasets.
\end{abstract}

\begin{IEEEkeywords}
heart disease prediction, classical machine learning, feature engineering, class imbalance, UCI Cleveland dataset
\end{IEEEkeywords}

\section{Introduction}
Heart disease prediction remains a widely studied tabular classification problem because it combines a clinically meaningful outcome with a compact set of structured predictors. The assignment provided for this mini-project requires a classical machine learning (ML) study, explicitly excludes deep learning methods, and identifies feature engineering and imbalance handling as primary improvement directions. The supplied dataset materials describe four heart disease databases, whereas both the archive notes and later published reviews indicate that the Cleveland subset is the standard basis for many machine learning experiments \cite{detrano1989cad,uciheart,ruqiya2023review,akella2021cad,rabbi2018performance}. The same archive materials also state that \texttt{processed.cleveland.data} is the reliable processed file and warn that \texttt{cleveland.data} is corrupted.

Within that context, the present work pursues three closely related objectives. First, it reproduces a benchmark-style experiment from the provided literature using standard classical models and the canonical binary labeling rule. Second, it evaluates an additional classical pipeline that incorporates feature engineering, class-balanced learning, and hyperparameter tuning. Third, it documents the resulting workflow in a structured IEEE-style paper accompanied by reproducible code and generated result files. The emphasis throughout is on verifiable reporting: statements about the dataset are drawn from the supplied archive files, while statements about model behavior are drawn from the generated outputs in the current workspace.

\section{Background and Related Work}
The University of California Irvine heart disease collection has been used for several decades as a benchmark for diagnostic classification. The original study cited in the archive describes probabilistic analysis for the diagnosis of coronary artery disease and reports approximately 77\% correct classification accuracy under the described experimental setting \cite{detrano1989cad}. The dataset repository further notes that published experiments commonly reduce the multiclass angiographic diagnosis to a binary formulation in which values greater than zero indicate disease presence \cite{uciheart}. This binary formulation is also reiterated in the review literature built around the Cleveland dataset \cite{ruqiya2023review,akella2021cad}. The present work adopts that same binary formulation because it is explicitly described in both the supplied dataset documentation and the published papers available in the workspace.

More recent studies continue to revisit this dataset family using standard supervised learning algorithms. Nasution \textit{et al.} compare Logistic Regression, Random Forest, Support Vector Machine, and K-Nearest Neighbors and report the highest accuracy from Random Forest, namely 89.7\%, without feature selection on a Kaggle-hosted UCI variant \cite{nasution2025heart}. Akella and Akella compare several machine learning algorithms on the Cleveland dataset and explicitly frame the study as a comparative analysis of multiple models \cite{akella2021cad}. Rabbi \textit{et al.} compare Support Vector Machine, K-Nearest Neighbors, and artificial neural networks under cross-validation on the Cleveland dataset \cite{rabbi2018performance}. Spencer \textit{et al.} evaluate feature-selection methods, including Chi-squared filtering, across multiple heart-disease datasets and classification algorithms \cite{spencer2020feature}. These studies are useful because they confirm that classical models remain competitive on this problem and provide a practical reference point for baseline reproduction.

At the same time, the reported literature does not present a fully uniform evaluation setting. Some studies combine multiple heart-disease datasets, some use single-site subsets, some report only a train--test split, and others rely on cross-validation. Feature-selection strategies also vary substantially, ranging from simple filter methods to more involved hybrid or optimization-based approaches. As a result, direct numerical comparison across papers should be treated cautiously. This is one reason the present work prioritizes an internally consistent experimental protocol on the supplied Cleveland file over attempting to claim exact comparability with every published result.

The supplied folder also contains papers that cite the dataset in other contexts, including privacy-preserving clinical decision systems \cite{ma2019ppcd}. Those papers demonstrate the continuing relevance of the dataset, but they are not direct benchmark papers for the current classification task. For that reason, Nasution \textit{et al.} \cite{nasution2025heart} was selected as the primary benchmark reference for reproduction. This choice is based on alignment with the assignment requirements rather than on a claim that it is the only possible or definitive benchmark.

From a methodological perspective, the reviewed literature can be grouped into three themes. The first theme is direct classifier comparison on the Cleveland-style tabular setting, typically using algorithms such as Logistic Regression, Random Forest, Support Vector Machine, K-Nearest Neighbors, and neural-network-based methods \cite{nasution2025heart,akella2021cad,rabbi2018performance}. The second theme is representation refinement, where authors study feature selection or feature construction to reduce redundancy and improve prediction quality \cite{spencer2020feature,rehman2025mi}. The third theme is data balancing or sampling strategy design for medical datasets that are small or class-imbalanced \cite{albert2023oversampling,muhammad2021smote}. The present project is most closely positioned at the intersection of the first and second themes, with limited engagement with the third through the use of class-weighted estimators instead of synthetic oversampling.

This positioning is important because it clarifies what the project does and does not attempt. It does not propose a novel classifier, a new medical dataset, or a clinically deployed system. Instead, it performs a controlled comparison of established classical methods, adds a modest amount of structured feature engineering, and evaluates whether those additions produce consistent benefits once leakage-free validation is enforced. In a course-project setting, this narrower contribution is more defensible than making broad clinical or methodological claims that cannot be supported by the available data.

\section{Dataset and Experimental Setup}
\subsection{Dataset Selection}
The Cleveland processed file contains 303 instances and 13 predictive attributes, with the target \texttt{num} converted to a binary label where \texttt{num} $>$ 0 indicates disease. The generated dataset summary in the current workspace confirms a class distribution of 164 negative cases and 139 positive cases. Missing values are limited to two attributes: \texttt{ca} contains four missing entries and \texttt{thal} contains two. The remaining attributes are complete in the processed file used in this project.

This dataset is small enough that experimental outcomes may change noticeably across validation schemes. In particular, an 80/20 holdout split leaves only 61 patients in the test partition. Accordingly, holdout evaluation alone would not be sufficient for careful interpretation. For that reason, the study reports both holdout and cross-validated baseline results, and uses cross-validation for tuned model selection.

\subsection{Preprocessing}
Two preprocessing regimes were used.
\begin{itemize}
\item Benchmark reproduction: median imputation for numeric variables, mode imputation plus one-hot encoding for categorical variables, and standard scaling for models that benefit from normalized numeric ranges.
\item Evaluated additional pipeline: the same base preprocessing plus engineered features such as \texttt{age\_oldpeak}, \texttt{thalach\_age\_ratio}, \texttt{exercise\_risk}, \texttt{risk\_cp}, and \texttt{vessel\_thal\_combo}.
\end{itemize}

The preprocessing choices were implemented within scikit-learn pipelines so that imputation, scaling, encoding, feature selection, and classification are executed consistently within each experiment. This is important because it avoids manual preprocessing steps that could be applied differently across models.

The benchmark paper discusses Chi-square and Mutual Information feature selection, but does not specify the retained feature count. Therefore, one explicit implementation assumption was required for reproduction: this study used \texttt{SelectKBest} with the top eight transformed features for both feature-selection variants. This assumption is stated here so that the reproduction is transparent rather than implicit. The broader choice of examining Chi-square and Mutual Information as filtering strategies is also consistent with published work on heart-disease or related coronary-disease tabular settings \cite{spencer2020feature,rehman2025mi}.

Another design choice concerns missing-data treatment. Because the processed Cleveland file contains missing values only in \texttt{ca} and \texttt{thal}, a simple imputation strategy is preferable to more complex data-completion methods. Median imputation for numeric variables and most-frequent imputation for categorical variables were therefore adopted within each pipeline. This decision keeps the preprocessing consistent across models and avoids introducing an additional modeling layer that would be difficult to justify from such a small sample.

The selection of evaluation metrics was similarly motivated by the structure of the task. Accuracy alone is easy to interpret, but it may hide asymmetric errors between disease-positive and disease-negative cases. Balanced accuracy and F1 score therefore provide additional information about class-sensitive behavior, while receiver operating characteristic area under the curve captures ranking quality independently of a fixed threshold. Reporting these metrics together makes the analysis more robust than relying on a single scalar summary.

\subsection{Models}
The benchmark family includes four classical models:
\begin{itemize}
\item Logistic Regression
\item Random Forest
\item Support Vector Machine with radial basis function kernel
\item K-Nearest Neighbors
\end{itemize}

The evaluated additional pipeline uses class-balanced Logistic Regression, class-balanced Support Vector Machine, and class-balanced Random Forest with grid search under 5-fold stratified cross-validation. The top two tuned models by cross-validated F1 score are then combined in a soft-voting ensemble. This design was chosen to remain within classical machine learning while still addressing one of the assignment's stated priorities, namely imbalance-aware modeling. Published heart-disease studies frequently discuss oversampling or imbalance-aware strategies such as the Synthetic Minority Oversampling Technique (SMOTE), especially in medical tabular settings \cite{albert2023oversampling,muhammad2021smote,rehman2025mi}. In the present implementation, imbalance was addressed through class-weighted estimators rather than synthetic oversampling. This choice avoids the additional distributional assumptions introduced by synthetic data generation and keeps the pipeline closer to the original observed dataset.

\subsection{Validation and Metrics}
To stay comparable with the reference paper, baseline results are reported on an 80/20 stratified holdout split with random seed 42. In addition, all baseline configurations are evaluated with 5-fold stratified cross-validation. Model selection for the additional tuned pipeline is also based on 5-fold stratified cross-validation. Evaluation uses accuracy, balanced accuracy, precision, recall, F1 score, receiver operating characteristic area under the curve, and confusion matrices. The use of cross-validation is consistent with several published heart-disease studies that rely on repeated partitioning or 10-fold validation to stabilize evaluation on relatively small datasets \cite{rabbi2018performance,albert2023oversampling,rehman2025mi}.

\section{Method}
The overall workflow is:
\begin{enumerate}
\item Load \texttt{processed.cleveland.data} and binarize the target.
\item Separate numeric and categorical variables.
\item Impute missing values and apply encoding/scaling inside pipelines.
\item Reproduce baseline models under three settings: no feature selection, Chi-square, and Mutual Information.
\item Engineer clinically motivated interaction features and train class-balanced tuned models.
\item Select the best cross-validated models and build a soft-voting ensemble.
\item Save all tables, figures, and summary files for paper generation.
\end{enumerate}

The engineered features were intentionally compact and interpretable. For example, \texttt{exercise\_risk} combines exercise-induced angina with ST depression, while \texttt{vessel\_thal\_combo} combines fluoroscopy vessel count with thalassemia status. Additional ratio-style terms, such as \texttt{thalach\_age\_ratio} and \texttt{chol\_trestbps\_ratio}, were introduced to test whether simple interactions could improve separability for linear or margin-based models. These terms were introduced as experimental interaction features within a classical machine learning pipeline; they are not claimed to be medically validated biomarkers.

The implementation was organized so that the same dataset is used throughout, namely \texttt{processed.cleveland.data}. This choice avoids mixing variants from external sources. The script also saves generated figures, summary tables, and JSON reports so that the narrative in the paper can be traced back to actual experimental outputs. Those generated artifacts include holdout tables, baseline cross-validation tables, tuned-model cross-validation tables, feature rankings, and ROC plots.

In the final revision of the code, bootstrap confidence intervals were also added for the main holdout metrics of each evaluated model. These intervals do not replace external validation, but they provide a useful sense of metric uncertainty on the small holdout sample. In addition, an ablation study was introduced to separate the contributions of feature engineering, class weighting, and ensembling within the Logistic Regression centered pipeline.

The method selection process is therefore two-stage. First, the reproduced baselines establish what can be achieved using standard classical models under transparent preprocessing and validation rules. Second, the additional tuned pipeline tests whether more deliberate design choices improve stability or ranking quality. This structure is aligned with the assignment requirement to reproduce a benchmark and then evaluate a justified improvement direction. The justification in this project rests not on a claim of novelty, but on the practical relevance of feature engineering, class-aware training, and model-selection discipline in small medical tabular datasets.

\section{Results}
\subsection{Dataset Overview}
Fig.~\ref{fig:classbalance} shows the class balance and Fig.~\ref{fig:corr} shows the feature correlation structure. Most variables have no missing data, while \texttt{ca} and \texttt{thal} require limited imputation only. The exported feature ranking places \texttt{ca}, thal-related indicators, \texttt{vessel\_thal\_combo}, \texttt{exercise\_risk}, and several chest-pain or slope indicators near the top.

The class balance plot confirms that the binary outcome is not severely skewed, but the classes are not perfectly balanced either. This provides a reasonable justification for examining balanced accuracy and for testing class-balanced estimators in the tuned pipeline. The heatmap is used descriptively rather than inferentially: it helps summarize linear relationships among the encoded variables, but it does not by itself establish causal or clinically meaningful dependence.

\begin{figure}[!t]
\centering
\includegraphics[width=0.85\columnwidth]{class_balance.png}
\caption{Binary class distribution in the Cleveland processed dataset.}
\label{fig:classbalance}
\end{figure}

\begin{figure}[!t]
\centering
\includegraphics[width=\columnwidth]{correlation_heatmap.png}
\caption{Correlation heatmap of the predictive variables and binary target.}
\label{fig:corr}
\end{figure}

\subsection{Benchmark Reproduction}
Table~\ref{tab:baselineholdout} summarizes the reproduced holdout results. On this single split, the strongest raw baseline is KNN without feature selection at 91.80\% accuracy and 0.9123 F1. However, the most relevant comparison to the benchmark paper remains Random Forest without feature selection, because Nasution \textit{et al.} report Random Forest as their best model at 89.7\% accuracy \cite{nasution2025heart}. Under the present implementation and dataset file, the reproduced holdout Random Forest reaches 90.16\% accuracy and 0.9000 F1.

Table~\ref{tab:baselinecv} summarizes the 5-fold cross-validation results for the same baseline family, computed on the training partition only. Under this criterion, the highest mean F1 is obtained by Logistic Regression without feature selection.

\begin{table}[!t]
\caption{Selected baseline reproduction results on the 80/20 holdout split}
\label{tab:baselineholdout}
\centering
\footnotesize
\setlength{\tabcolsep}{4pt}
\resizebox{\columnwidth}{!}{
\begin{tabular}{lcccc}
\toprule
Model & Variant & Acc. & F1 & ROC-AUC \\
\midrule
KNN & None & 0.9180 & 0.9123 & 0.9627 \\
Random Forest & None & 0.9016 & 0.9000 & 0.9394 \\
SVM & None & 0.9016 & 0.8966 & 0.9610 \\
Logistic Regression & None & 0.8689 & 0.8667 & 0.9578 \\
Logistic Regression & Mutual Info & 0.8852 & 0.8814 & 0.9448 \\
Random Forest & Chi-square & 0.8689 & 0.8621 & 0.9448 \\
\bottomrule
\end{tabular}
}
\end{table}

\begin{table}[!t]
\caption{Selected baseline 5-fold cross-validation results}
\label{tab:baselinecv}
\centering
\footnotesize
\setlength{\tabcolsep}{4pt}
\resizebox{\columnwidth}{!}{
\begin{tabular}{lcccc}
\toprule
Model & Variant & Mean Acc. & Mean F1 & Mean ROC-AUC \\
\midrule
Logistic Regression & None & 0.8470 & 0.8254 & 0.9065 \\
SVM & Mutual Info & 0.8347 & 0.8127 & 0.8878 \\
SVM & Chi-square & 0.8304 & 0.8105 & 0.8888 \\
Logistic Regression & Chi-square & 0.8261 & 0.8080 & 0.9059 \\
SVM & None & 0.8305 & 0.8041 & 0.8850 \\
KNN & Chi-square & 0.8261 & 0.8017 & 0.8858 \\
\bottomrule
\end{tabular}
}
\end{table}

\noindent The holdout and cross-validated rankings are not identical. This is a meaningful result rather than a contradiction. The holdout split captures one particular partition of a small dataset, whereas 5-fold cross-validation averages performance across multiple partitions. Therefore, the tables are reported together so that the reader can distinguish split-specific performance from average validation behavior.

An additional point follows from the baseline tables. The strongest single holdout result is produced by K-Nearest Neighbors, whereas the strongest mean cross-validated baseline result is produced by Logistic Regression. This pattern implies that the dataset does not favor one family of models under every evaluation criterion. Distance-based models appear capable of fitting this particular split very effectively, whereas the simpler linear baseline appears more stable on average across the training-fold partitions. That contrast is one of the central empirical observations of the study.

\subsection{Additional Tuned Pipeline}
Cross-validation on the training partition selected balanced Logistic Regression as the top candidate by mean F1 score, followed by balanced SVM. Their tuned settings were $C=0.5$ with $L_1$ penalty for Logistic Regression and $C=0.5$ with \texttt{gamma=scale} for SVM. The resulting soft-voting ensemble achieves 88.52\% accuracy, 0.8814 F1, balanced accuracy of 0.8885, and 0.9600 ROC-AUC on the holdout split.

Table~\ref{tab:improved} shows the evaluated-model results. Compared with the benchmark paper's reported best Random Forest accuracy of 89.7\% \cite{nasution2025heart}, the evaluated ensemble is slightly below that number on the present split. Compared with the reproduced holdout Random Forest baseline in this study, the ensemble has higher ROC-AUC but lower accuracy and F1. Compared with the strongest single-split raw baseline in this experiment, the ensemble is clearly weaker in accuracy and F1. In addition, the holdout metrics of the ensemble and the tuned balanced Logistic Regression model are identical for accuracy, balanced accuracy, precision, recall, and F1, which means the ensemble does not provide a clear empirical advantage in those metrics on this split.

This result is important because it prevents an overly optimistic interpretation of the additional pipeline. The tuned approach remains competitive and exhibits high holdout ROC-AUC, but its main value lies in disciplined evaluation and clearer tradeoff analysis rather than in a large absolute gain over the best simple baseline. In other words, the project's strongest message is methodological rather than algorithmic.

\begin{table}[!t]
\caption{Training-partition cross-validated tuned models and holdout results}
\label{tab:improved}
\centering
\footnotesize
\setlength{\tabcolsep}{4pt}
\resizebox{\columnwidth}{!}{
\begin{tabular}{lccc}
\toprule
Model & Acc. & F1 & ROC-AUC \\
\midrule
Balanced Logistic Regression & 0.8852 & 0.8814 & 0.9643 \\
Balanced SVM & 0.8689 & 0.8667 & 0.9675 \\
Soft Voting Ensemble & 0.8852 & 0.8814 & 0.9600 \\
\bottomrule
\end{tabular}
}
\end{table}

\begin{table*}[!t]
\caption{Ablation study for the Logistic Regression centered pipeline}
\label{tab:ablation}
\centering
\footnotesize
\setlength{\tabcolsep}{5pt}
\begin{tabular}{lcccc}
\toprule
Variant & Holdout Acc. & Holdout F1 & Mean CV F1 & Mean CV ROC-AUC \\
\midrule
\makecell[l]{Baseline\\Logistic Regression} & 0.8689 & 0.8667 & 0.8254 & 0.9065 \\
\makecell[l]{Feature\\Engineering Only} & 0.8689 & 0.8667 & 0.8156 & 0.9047 \\
\makecell[l]{Class\\Weights Only} & 0.8689 & 0.8667 & 0.8397 & 0.9097 \\
\makecell[l]{Feature Eng.\\+ Class Weights} & 0.8852 & 0.8814 & 0.8335 & 0.9085 \\
\makecell[l]{Soft Voting\\Ensemble} & 0.8852 & 0.8814 & 0.8131 & 0.8997 \\
\bottomrule
\end{tabular}
\end{table*}

\begin{figure}[!t]
\centering
\includegraphics[width=0.92\columnwidth]{roc_curves.png}
\caption{ROC curves for the tuned proposed models on the holdout split.}
\label{fig:roc}
\end{figure}

\begin{figure}[!t]
\centering
\includegraphics[width=\columnwidth]{top_features.png}
\caption{Top-ranked transformed features from the tuned balanced Logistic Regression model.}
\label{fig:topfeatures}
\end{figure}

\subsection{Interpretation}
The exported feature ranking places \texttt{ca}, thal-related indicators, the interaction \texttt{vessel\_thal\_combo}, and \texttt{exercise\_risk} near the top. This is an observation from the trained models in this project. It should not be interpreted as a clinical conclusion without external validation.

Similarly, the ROC curves indicate that the tuned models and ensemble maintain strong ranking behavior on the holdout split. This statement is restricted to the current dataset split and generated ROC-AUC values. It does not establish superiority on external datasets, because no external validation data were available in the supplied materials. Notably, the tuned SVM attains the largest holdout ROC-AUC among the tuned models even though its holdout accuracy and F1 are lower than those of the tuned balanced Logistic Regression model.

The bootstrap intervals further reinforce the need for caution. For the final ensemble, the bootstrap mean holdout accuracy is approximately 0.8851 with a 95\% interval from 0.8033 to 0.9508, while the bootstrap mean holdout F1 is approximately 0.8797 with a 95\% interval from 0.7777 to 0.9565. These intervals are wide relative to the differences among several competing models, which indicates that small metric gaps should not be overinterpreted on this sample size.

\subsection{Error Analysis}
The corrected holdout confusion matrix of the final ensemble contains 28 true negatives, 26 true positives, 5 false positives, and 2 false negatives. This means the evaluated tuned pipeline misses relatively few positive cases on this split, but it produces more false alarms than the strongest reproduced KNN baseline. In a clinical decision-support setting, that tradeoff may or may not be preferable depending on the cost of missed disease versus unnecessary follow-up testing. The present study does not assign a clinical utility function to those outcomes, so the confusion-matrix interpretation is reported descriptively only.

The class-specific report further shows that the ensemble's recall for the positive class is 0.9286, while recall for the negative class is 0.8485. This asymmetry is consistent with the use of class-balanced estimators, which push the model toward better sensitivity for the disease-present class. However, because the dataset is small and no external validation cohort is available, this behavior should not be generalized beyond the current split.

A second observation concerns the relation between ROC-AUC and threshold-based metrics. The tuned Support Vector Machine attains the strongest holdout ROC-AUC among the tuned models, yet it does not achieve the best holdout F1 score. This illustrates a familiar but important distinction: ranking quality and thresholded classification quality are related but not identical. In practical terms, a model can rank positive cases well while still producing a less favorable fixed-threshold confusion matrix. This distinction supports the decision to report multiple evaluation metrics rather than a single headline accuracy value.

\subsection{Uncertainty and Stability}
One of the most important findings of this project is that the ordering of models changes across evaluation schemes. KNN attains the strongest single-split holdout performance, Logistic Regression attains the strongest mean baseline cross-validated F1, and the tuned SVM attains the highest holdout ROC-AUC among the tuned models. The existence of these different rankings is itself evidence that the dataset is too small to support broad superiority claims from a single metric alone.

The bootstrap confidence intervals support the same conclusion. For the final ensemble, the 95\% interval for holdout ROC-AUC ranges from 0.9000 to 0.9989, and the 95\% interval for holdout F1 ranges from 0.7777 to 0.9565. These ranges are sufficiently wide that several of the competing models remain plausibly similar in practical performance. Accordingly, the uncertainty analysis strengthens the paper's central conclusion that rigor in evaluation is more defensible than claiming a decisive model improvement.

\subsection{Practical Implications}
Although the present study is not a deployment paper, the observed tradeoffs are still informative from a practical perspective. If the design objective is to maximize simple holdout accuracy on the current split, the reproduced K-Nearest Neighbors baseline is the strongest choice. If the objective is to favor a simpler and more stable model family under cross-validation, Logistic Regression is a reasonable candidate. If the objective is to retain strong ranking performance while biasing the system toward higher sensitivity for disease-positive cases, the tuned class-balanced models remain relevant. These observations suggest that model selection in small medical tabular problems should be aligned with the intended operating criterion rather than with a single generic metric.

The experiments also show that modest methodological choices can substantially affect reported conclusions. A result based on a single holdout split may favor one model, whereas a training-partition cross-validation summary may favor another. Likewise, class weighting can improve average F1 score without guaranteeing the strongest holdout result. From an engineering viewpoint, this means that reproducibility, leakage control, and metric selection are not secondary implementation details; they directly affect the scientific interpretation of the project.

\section{Discussion}
Several lessons emerged from the experiments.

First, the provided benchmark paper is reproducible only approximately from the available information. It uses a Kaggle-hosted UCI variant, while the assignment provides the original archive files. It also omits some implementation details, such as the exact feature-selection configuration. This project therefore reports an explicit approximation rather than an exact replication claim. That limitation is not unusual in this literature, where reported experimental settings sometimes differ across dataset variants, preprocessing decisions, and validation protocols.

Second, small datasets can produce noticeable metric changes across evaluation setups. On the chosen holdout split, untuned KNN without feature selection produces the best accuracy and F1. Under 5-fold baseline cross-validation on the training partition, Logistic Regression without feature selection attains the highest mean F1 among the reproduced baselines. Cross-validation for the tuned pipeline also favors balanced Logistic Regression, while the final ensemble does not improve holdout F1 beyond that tuned Logistic Regression model. This demonstrates that a single performance ranking should not be treated as universal when the sample size is limited.

Third, the ablation results provide a more precise interpretation of the added pipeline. In this project, class weighting contributes more to cross-validated improvement than feature engineering alone, while the soft-voting ensemble does not improve the mean cross-validated F1 over the stronger single tuned Logistic Regression variant. The tuned feature-engineered Logistic Regression still improves the holdout metrics relative to plain Logistic Regression, but the gain is modest and remains below the strongest simple KNN holdout baseline. From a strict reviewer perspective, the main contribution of the tuned pipeline in its current form is methodological discipline and controlled comparison rather than a decisive empirical improvement. The literature also suggests that preprocessing, feature filtering, and data-balancing choices can materially affect reported outcomes on heart-disease prediction benchmarks \cite{spencer2020feature,albert2023oversampling,rehman2025mi}.

Fourth, the project highlights the difference between selecting a model for reporting and selecting a model for deployment. If the task were purely to maximize holdout accuracy on this split, K-Nearest Neighbors would be the most attractive model. If the task were to favor stable cross-validated baseline performance on the training partition, Logistic Regression would be more attractive. If the task prioritized strong ranking behavior with sensitivity to positive cases, the tuned models and ensemble remain relevant. This ambiguity is precisely why small medical benchmarks demand careful justification rather than a simple leaderboard mentality.

Fifth, the project underscores the value of negative or mixed findings in empirical machine learning. In a small study such as this one, it would be easy to present only the most favorable result and describe it as an improvement. However, a stricter academic standard requires the opposite approach: all relevant comparisons must be shown, and the final interpretation should reflect the full evidence rather than the most convenient subset of it. The present paper therefore treats the inability of the tuned pipeline to surpass the strongest simple baseline not as a failure of the study, but as a meaningful outcome that clarifies what the data do and do not support.

There are also limitations. The study uses only 303 instances from a single processed file. There is no external validation cohort, no temporal split, and no patient-level metadata beyond the supplied tabular attributes. Although the current implementation now avoids holdout leakage by restricting cross-validation and tuning to the training partition, the dataset remains too small to support strong claims about generalization. The bootstrap intervals and variation between holdout and cross-validation rankings both support that conclusion. Furthermore, the paper template could not be compiled in the current environment because a local \texttt{pdflatex} installation is unavailable. Consequently, the final page count and compiled formatting cannot be verified within this workspace. Future work may evaluate the pipeline on multi-center combinations of the UCI sites, compare threshold-tuning strategies, report calibration error in addition to discrimination metrics, and assess whether the observed model rankings remain stable under repeated resampling.

Several concrete extensions are especially well motivated by the current findings. First, repeated or nested cross-validation would provide a stronger estimate of model stability than a single 5-fold split. Second, a threshold-tuning analysis could determine whether the tuned models can trade a small amount of specificity for improved disease recall in a more controlled manner. Third, calibration analysis could reveal whether the probability outputs are reliable enough for decision support rather than only for ranking. Fourth, integrating multiple UCI heart-disease subsets with a carefully harmonized preprocessing protocol may help determine whether the conclusions observed on the Cleveland file persist in a broader multi-source setting. Each of these directions would strengthen the study without requiring unsupported claims in the current paper.

\subsection{Reproducibility}
The code accompanying this paper stores the experiment logic in a single executable script and exports intermediate artifacts, final tables, and summary JSON files to the \texttt{results} directory. A dependency file is also included so that the Python environment can be reconstructed more directly. This does not make the project fully archival in the sense of a large benchmark release, but it substantially improves traceability compared with manually copied experimental outputs.

An additional practical strength is that the final manuscript claims can be traced back to concrete generated artifacts. The holdout comparison tables, cross-validation summaries, bootstrap intervals, and ablation outputs are all exported from the same script rather than copied manually from ad hoc experiments. This reduces the risk of internal inconsistency between the code and the paper narrative, which is a common weakness in small research projects.

\section{Conclusion}
This project delivered a classical machine learning workflow for the provided heart disease task, including code, result files, figures, and an IEEE-style paper draft. After reading the supplied archive contents, \texttt{processed.cleveland.data} was selected as the experimental source. Baseline reproduction showed strong performance from simple classical methods, but the specific model ranking depended on whether the evaluation used a single holdout split or 5-fold cross-validation on the training partition. The evaluated tuned class-balanced pipeline did not surpass the strongest simple baseline on the corrected holdout evaluation, although it remained competitive and achieved strong ROC-AUC values. The ablation study further showed that class weighting was more beneficial than feature engineering alone within the Logistic Regression centered design, and the uncertainty analysis showed that small metric differences should be interpreted cautiously. The principal conclusion supported by the verified results is therefore not that a more elaborate pipeline is universally superior, but that careful validation, leakage-free evaluation, uncertainty reporting, and explicit methodology are essential when drawing conclusions from small medical tabular datasets.

At the same time, the study still satisfies the assignment's broader educational objective. It demonstrates how to move from dataset documentation and benchmark reading to a reproducible experimental workflow, how to test plausible improvement directions without overstating the outcome, and how to present negative or mixed findings in an academically responsible way. In that sense, the final paper supports a more mature conclusion: a well-executed baseline with honest evaluation can be more valuable than a superficially stronger but poorly validated improvement claim.

\section*{Acknowledgment}
This report should acknowledge the original UCI Heart Disease data contributors named in the dataset archive: Andras Janosi, William Steinbrunn, Matthias Pfisterer, and Robert Detrano.

\bibliographystyle{IEEEtran}
\bibliography{references}

\end{document}
